\section{Degree-separated coding: a ten-integer equation}
\label{sec:degree-separated-sun}

There is a completely fixed integer polynomial
\[
 F_{10}(f,g,h,k,l,w,x,y,z,m)
\]
of exact total degree $8\,469\,528$, evaluated by a division-free circuit of
$533$ binary arithmetic operations, for which ordinary mathematics
formalizable in ZFC proves
\[
 \neg\exists(f,g,h,k,l,w,x,y,z,m)\in\mathbb Z^{10}\;(F_{10}=0)
 \quad\Longleftrightarrow\quad\operatorname{Con}(\mathrm{ZFC}).
\]
This is a new source-coding architecture, not a new unknown-count theorem.
The inherited compiler calibration and Sun arithmetic bridge remain stated
mathematical dependencies. No consistency decision, global minimality claim,
or new end-to-end proof-assistant certificate is asserted.

The structural change is to split the extraction convolution by homogeneous
degree. A linear label may now collide with a quadratic label: their kernels
are separate summands. This relaxes the finite isolation problem rather than
merely finding a better placement for the previous problem. Elementary row
operations, one shared positivity shift, borrow-tolerant target padding,
and exact mask-population calibration complete the construction.

\subsection{Fixed source and three homogeneous kernels}

Use exactly the unshifted, strictly positive, 31-coordinate parent in
\texttt{data/parent\_quartic31.json}, with its complete original four-port
binding $(\beta_{\rm native},K_0,M_C,M_F)$. The four included compressed
payloads have bit lengths $(1064921,1070004,1064921,1064922)$; the last is
\texttt{MF\_source}, already shifted, and is not shifted again.
\texttt{INSTANCE.json} is the active binding. No package-52 transport row or
coefficient is spliced into this tuple. The pinned parent theorem is
\[
 \exists p\in\mathbb Z_{>0}^{31}\;r_1(p)=\cdots=r_{22}(p)=0
 \quad\Longleftrightarrow\quad\neg\operatorname{Con}(\mathrm{ZFC}).
\]
Replace only $r_{14}$ by $\rho_{14}=r_{14}+r_2$ and $r_{15}$ by
$\rho_{15}=r_{15}+r_3$; otherwise put $\rho_j=r_j$. The inverse row
operations preserve the common zero set over every coefficient tuple.
There are now 89 records: seven constant, 39 linear, and 43 quadratic.

In the parent coordinate order
\[
\begin{split}
 &(J,F,\alpha,z,f,h,i,o,j,s,w,\tau,\eta,\zeta,y,\gamma,\\
 &\hspace{22mm}N,E,P,c,a,X,D,C_2,L,F_2,V,g,I,B_{\rm aux},\Delta),
\end{split}
\]
choose the positive, distinct weights
\[
\begin{split}
 (w_1,\ldots,w_{31})={}&(1,7,61,36,16,85,20,67,43,11,34,13,96,80,29,50,\\
 &2,109,110,112,116,91,53,55,27,101,92,9,107,37,60).
\end{split}
\]
The row order and targets are
\[
\begin{split}
 &(2,22,8,15,9,11,3,17,13,18,12,20,4,21,19,1,6,10,16,14,7,5),\\
 &(28,85,286,384,593,602,642,731,906,966,1083,1212,1231,1393,\\
 &\hspace{16mm}1461,1471,1678,1719,1732,1880,2107,2129).
\end{split}
\]
Let $q_j$ denote the target assigned to row $j$, and write
$\rho_j=\sum_\alpha a_{j,\alpha}p^\alpha$. For $d=|\alpha|$ put
\[
 e(\alpha)=\sum_i\alpha_iw_i,\qquad
 c_{j,\alpha}=
 \begin{cases}2a_{j,\alpha},&d=0,1,\\
 a_{j,\alpha}\prod_i\alpha_i!,&d=2.
 \end{cases}
\]
Define the fixed integer polynomials
\[
 D_d(Z)=\sum_{j,\,|\alpha|=d}c_{j,\alpha}Z^{q_j-e(\alpha)},\qquad
 U_p(Z)=\sum_i p_iZ^{w_i},\qquad
 G_p(Z)=D_2(Z)U_p(Z)^2+D_1(Z)U_p(Z)+D_0(Z).
\]
All reversed exponents are nonnegative. The exact finite certificate checks,
for every used degree-$d$ monomial $\alpha$ and every degree-$d$ monomial
$\gamma$ in 31 variables,
\[
 q_j-e(\alpha)+e(\gamma)=q_k
 \quad\Longleftrightarrow\quad j=k\text{ and }\alpha=\gamma.
\]
There are respectively $1,31,496$ expansion monomials in degrees $0,1,2$,
and 22,544 source/expansion interactions. There is no requirement to
separate a degree-one label from a degree-two label. Multinomial expansion
and the displayed coefficient weights therefore give the universal identities
\begin{equation}\label{eq:sun-split-extraction}
 [Z^{q_j}]G_p(Z)=2\rho_j(p)\quad(1\le j\le22).
\end{equation}
The independent checker compares the entire extracted polynomials in all
31 coordinates and four ports, not just evaluations at test points.

Write $M=116$, $Q=2129$. The exact kernel degrees and convolution height are
\[
 \deg D_0=1719,\qquad\deg D_1=2027,\qquad\deg D_2=1929,\qquad
 H=\max_d(\deg D_d+dM)=2161.
\]
The leading term of $D_2$ is exactly $+Z^{1929}$, arising from $EX$ in
row five. Every kernel separately has distinct exponents. Replacing each
port monomial in a coefficient by its strict bit bound and taking absolute
integer coefficients gives a maximum majorant of bit length $2134927$.
Consequently, with
\[
 C_0=2^{2134927},\qquad\kappa=2134940,\qquad\beta_S=2^\kappa,
\]
every kernel coefficient has absolute value less than $C_0$, and
$\beta_S>4096C_0$.

\subsection{Source polynomials and pre-binary bounds}

In the unrestricted integer unknowns $f,g$ define
\begin{gather*}
 b=1+3f,\quad B=\beta_S b^2,\quad\ell=\sum_i B^{w_i},\quad U_c=\ell+g,\\
 G_c=D_2(B)U_c^2+D_1(B)U_c+D_0(B),\quad A_T=\sum_j B^{q_j},\\
 N_0=B^{117},\quad N_1=B^{2162},\quad N=N_0N_1=B^{2279},\\
 M_0=N_0-1-(b-1)\ell,\quad J_c=2G_c+A_T,\\
 S_c=g+J_cN_0,\quad T_c=M_0+(B-2)A_TN_0,\\
 X=N^2+(S_c+T_c)N+N-1-S_c,\\
 Y=\beta_S^{15}\,2^{\kappa-21}B^{2402}b
   =2^{2418\kappa-21}b^{4805},\qquad
 \mathcal C=bB^{116}=2^{116\kappa}b^{233}.
\end{gather*}
All powers have fixed nonnegative integer exponents. Every displayed
register is an abbreviation, never an additional quantified unknown.

Assume $f\ne0$, $b$ is a square, and $0\le g<2\mathcal C$.
Then $b\equiv1\pmod3$ implies $b\ge4$, and $B>4096C_0b^2$.
Distinct positive weights give $\ell<2B^{116}$, hence
\[
 B^{116}\le U_c<3bB^{116}.
\]
In particular, without assuming $b$ is a power of two,
\[
 D_2(B)>\tfrac12B^{1929},\quad
 |D_1(B)|<2C_0B^{2027},\quad |D_0(B)|<2C_0B^{1719}.
\]
After division by $B^{1929}U_c^2$, the two possible negative terms
are bounded by $2C_0/B^{18}<1/8$ and $2C_0/B^{442}<1/8$.
Thus $G_c>\tfrac14B^{1929}U_c^2>0$. This is the necessary signed
leading-layer argument; positivity is not inferred from positive inputs
alone.

The kernel degree bounds also give
\[
 |G_c|<2C_0(1+3b+9b^2)B^H\le26C_0b^2B^H,\qquad
 2G_c<N_1/64.
\]
Since $Q=2129<H$, both $A_T$ and $(B-2)A_T$ are less than $N_1/64$.
It follows that
\[
 0<J_c<N_1/32,\qquad J_c+(B-2)A_T+2<N_1.
\]
The digits of $M_0$ are $B-1$, except that the 31 designated cells
have digit $B-b$. Therefore $0\le M_0<N_0$, and $g<N_0$ as well.
The preceding strict high-block bound proves
\begin{equation}\label{eq:split-tight-range}
 0\le S_c,\quad0\le T_c,\qquad S_c+T_c<N.
\end{equation}
Consequently $X\ge N^2>3b$ and $Y>\max(b,256)$. Because $\kappa$ is
even, $B\equiv N\equiv1\pmod3$, $M_0\equiv0$, and
$T_c\equiv-22$; hence $X\equiv1+T_c\equiv0\pmod3$.
These are all established before binary-base recovery.

\subsection{Borrow-tolerant extraction and exact mask population}

Suppose now that $b$ is a power of two. So are $B,N_0,N$. The low
mask condition $g\mathbin{\&}M_0=0$ is exactly
\[
 g=\sum_i n_iB^{w_i},\quad0\le n_i<b.
\]
For these digits put $p_i=n_i+1$, so $U_c=\sum_i p_iB^{w_i}$ encodes
positive parent coordinates without expanding a successor substitution
in every residual. Write $h_s=[Z^s]G_p(Z)$ and let $\epsilon_s$ be the
indicator of the 22 targets. All coefficients, including unsupported
collisions, satisfy
\[
 |h_s|<C_0\{(\sum_i p_i)^2+\sum_i p_i+1\}
 \le993C_0b^2<B/4.
\]
The signed expansion of $J_c$ is $\sum_{s=0}^{H}(2h_s+\epsilon_s)B^s$.
Euclidean digit normalization has incoming carries only in $\{-1,0\}$:
\[
 u_0=0,\quad u_{s+1}=\left\lfloor\frac{2h_s+\epsilon_s+u_s}{B}\right\rfloor,
 \quad d_s=2h_s+\epsilon_s+u_s-Bu_{s+1}.
\]
At a marked cell the mask $B-2$ admits precisely digits $0,1$.
If $h_s=0$, the unnormalized digit $1+u_s$ passes. If $h_s\ge1$,
it lies in $[2,B)$ and fails. If $h_s\le-1$, it lies strictly between
$-B+1$ and zero, and its normalized digit is at least two, so it fails.
Thus incoming borrows do not affect the zero test. Positivity of $J_c$
was proved independently, ensuring these are its ordinary digits.
Combining the disjoint low and high blocks with
\eqref{eq:sun-split-extraction} proves
\[
 S_c\mathbin{\&}T_c=0
 \quad\Longleftrightarrow\quad
 g=\sum_i n_iB^{w_i},\ 0\le n_i<b,\quad r_1(n+1)=\cdots=r_{22}(n+1)=0.
\]

For $N=2^u$ and $S,T\ge0$ with $S+T<N$, let
$X_*=N^2+(S+T)N+N-1-S$. Its disjoint binary blocks give
\[
 s_2(X_*)=1+u+s_2(T)-\{s_2(S)+s_2(T)-s_2(S+T)\}.
\]
The braces count carries with multiplicity, and vanish precisely when
$S\mathbin{\&}T=0$. Legendre's formula therefore implies
\[
 2^{1+u+s_2(T)}\mid\binom{2X_*}{X_*}
 \quad\Longleftrightarrow\quad S\mathbin{\&}T=0.
\]
The stronger range \eqref{eq:split-tight-range} is essential.
Writing $b=2^k$, $B=2^v$, the actual mask population is
\[
 s_2(T_c)=(117+22)v-31k-22,\qquad u=2279v.
\]
Hence $1+u+s_2(T_c)=2418v-31k-21=\log_2Y$, exactly explaining the
fixed polynomial threshold used above.

Conversely, a positive parent zero is not the all-one vector: original
row five equals one there. For arbitrarily large $b=4^s$ exceeding all
its coordinates, encode $n_i=p_i-1$. Then $g\ge B>b$ and $g<\mathcal C$;
the latter follows from $(b-1)\ell<bB^{116}$. All marked tests pass and
$Y\mid\binom{2X}{X}$. This proves both directions of the entire source
contract needed by the inherited finalizer, including its lower bound on
forward witnesses and its pre-binary reverse bounds.

\subsection{Explicit radical-free finalizer}

The following formulas are the inherited package-31 quintic,
merged-square, structural-weight finalizer. They are written out to fix
the actual equation. Distinguish its local $B_*,S_*,T_*$ from the source
base and masks above. Define
\begin{gather*}
 L=lY,\qquad U=2LX,\qquad V=4gwY,\qquad W=bw,\\
 A=U(V+1),\qquad B_*=2X+1,\qquad C=B_*+(A-2)h,\\
 P=U^2V,\qquad K=X+1+k(P-2),\\
 S_*=2A-5,\qquad T_*=3WC-2(W^2-1),\\
 R_2=(P^2-4)K^2+4,\\
 \Delta=8\mathcal C^3gK^2-g^2\{32(C-KL)^2\mathcal C^3+g^2K^2\},\\
 O=f^2l^2\Delta,\qquad \xi=zx,\\
 D=(A^2-4)C^2+4,\qquad E=C^2D\xi,\\
 F=4(A^2-4)E^2+1,\qquad G_*=2(A+2)(A-2)^2E^2-1,\\
 \Omega=16E,\qquad H_*=C+(B_*-C)F+\Omega Fy,\\
 I_*=(G_*^2-1)H_*^2+1.
\end{gather*}
Put
\begin{align*}
 r&=T_*+S_*y,&\Sigma&=\xi S_*^5,\\
 \Theta&=(r^3-rS_*^2-S_*^3)(r^2+23S_*^2),\\
 a_0&=b\{1+(3O-4)b^2z^2\},&a_1&=R_2,&a_2&=DFI_*,\\
 w_0&=1,&w_1&=a_0,&w_2&=a_0(P^2+K^2+3),\\
 Z&=\Sigma m-\Theta,&A_i&=\Sigma^2w_i^2a_i\quad(0\le i\le2),\\
 V_n&=Z^2+A_2-A_0-A_1,&W_n&=V_n^2+4Z^2A_2-4A_0A_1.
\end{align*}
The promised single equation is
\begin{equation}\label{eq:split-final-ten}
 \boxed{\ F_{10}=W_n^2-16V_n^2Z^2A_2=0.\ }
\end{equation}
Every expression is an integer polynomial. In particular, radicals and
binomial coefficients appear in the proof only, never in the equation.

\paragraph{Exact dependency boundary.}
Sun's 2021 Theorems 4.12--4.13, specialized as in his 2026 Propositions
6.1--6.2, provide the two bridge directions for any source satisfying the
contract just proved. Forward witnesses have $b$ square, $R_2$ square,
$S_*\mid T_*$, $C=u_{B_*}(A)$ and $O>0$. Reverse witnesses with
$b$ square, $O>0$, $\xi\ne0$, $R_2,DFI_*$ square and $S_*\mid T_*$
recover $b\in2\uparrow$, $0<g<2\mathcal C$ and the required binomial
divisibility. The residue interpolation theorem permits $\xi$ to be
divisible by any prescribed positive integer and permits any residue
class for $y\pmod\xi$.
The only data substituted here are the newly proved source polynomials;
the bridge does not require $Y=N^2$.

For completeness, the inherited finalizer implications can be verified
as follows. The polynomial $(v^3-v-1)(v^2+23)$ has a root modulo every
positive integer and none over $\mathbb Q$. For odd primes other than
23 the cubic discriminant is $-23$: if the cubic is irreducible its
Frobenius permutation is even, so the quadratic has a simple root;
otherwise the cubic has a simple root. At 23 the cubic has simple root
3. At 2, $-23\equiv1\pmod8$. Hensel lifting and the Chinese remainder
theorem finish the local-root argument. Prime valuations in the two
homogeneous factors show that $\Sigma\mid\Theta$ forces $S_*\mid T_*$:
if $v_p(S_*)>v_p(T_*)$, the valuation of $\Theta$ is
$5v_p(T_*)<5v_p(S_*)$. Since $S_*$ is odd, it never vanishes; the absence
of rational roots ensures $\Theta\ne0$.

When $z\ne0$, $a_0$ square implies that its coprime factors $b$ and
$1+(3O-4)b^2z^2$ are signed squares. The congruence $b\equiv1\pmod3$
excludes a negative square and excludes zero, including the zero-factor
exception. Thus $b$ is a positive square and the usual positivity Pell
lemma gives $O>0$. Conversely, forward witnesses have $O\ge4$: the
positive integer $\Delta$ is divisible by $g\ge b\ge4$. For
$d=3O-4>0$, $d\equiv2\pmod3$ is nonsquare. A positive Pell unit raised
to its order modulo $b$ gives $v=bz>0$ with $u^2-dv^2=1$, making $a_0$
a square. Residue interpolation then supplies $\xi=zx$ and a $y$ that
simultaneously makes $a_2$ square and $\Sigma\mid\Theta$.

If $z\ne0$, the displayed $a_0$ always has $|a_0|>1$. The exceptional
possibilities $a_0=0,\pm1$ force $|b|=1$, hence $b=1$, $f=O=0$, and
$a_0=1-4z^2$, a contradiction. Also
$\sqrt{|a_1|}\le P^2+K^2+2$. Thus the weights obey
\[
 |w_1|>\sqrt{|a_0|},\qquad
 |w_2|>\sqrt{|a_0|}+|w_1|\sqrt{|a_1|}.
\]
In the multiquadratic field, an automorphism changing a nonzero weighted
root has a largest changed index; its magnitude exceeds the sum of all
smaller changed magnitudes. A rational weighted sum can therefore occur
only when every $a_i$ is an integer square. This also covers negative
radicands by complex absolute values, and zero radicands separately.

The exact norm identity
\[
 F_{10}=\prod_{\epsilon\in\{\pm1\}^3}
 \left(Z+\sum_{i=0}^2\epsilon_i\sqrt{A_i}\right)
\]
has 35 expanded terms and is checked independently by integer polynomial
arithmetic. If $\Sigma=0$ then $F_{10}=\Theta^8\ne0$. At a zero of
$F_{10}$ we therefore have $\Sigma\ne0$, hence $z\ne0$; separation forces
all $a_i$ square and $\Sigma\mid\Theta$. The reverse bridge and source
proof recover a parent solution. Conversely the forward construction
just described makes all $a_i$ square and $\Sigma\mid\Theta$; choosing
$m$ to cancel one signed factor gives $F_{10}=0$. Composition with the
pinned parent theorem proves the stated consistency equivalence in ZFC.


\subsection{Exact circuit, degree, and reproducibility}

The actual circuit uses the same finalizer polynomial in factored form.
For example, with $A=2A_{1/2}$ and $F-1=4(A^2-4)E^2$, its Lucas expression
is $G_*=(A_{1/2}-1)(F-1)-1$; also
$H_*=C+F(B_*-C+\Omega y)$. The common factor
$\Sigma^2a_0^2$ is reused in the last two scaled radicands.
Fixed powers reuse already computed powers $B^i,B^j$ when $i+j$ is the
required exponent, falling back to binary exponentiation otherwise.
These are exact polynomial identities, not a relaxation of the equation.

\texttt{generated/sun10.json} contains all ten variable leaves and 533
live binary gates: 340 multiplications, 166 additions, and 27 subtractions.
Multiplication by a numeral is charged. The fixed coefficient leaves are
specified by \texttt{INSTANCE.json}, the four payloads, and
\texttt{generated/source\_records.json}. Their construction is not included
in this arithmetic-operation count, which is not bit complexity.

The source degrees are $\deg X=9116$, $\deg Y=4805$, and
$\deg\mathcal C=233$. The propagated final upper bound is
\[
 456(9116)+896(4805)+24(233)+1760=8\,469\,528.
\]
After all ten variables are replaced by one indeterminate $T$, the
coefficient of $T^{8469528}$ has nonzero residues
\[
 \begin{array}{c|ccc}
 p&101&1009&1000003\\\hline
 [T^{8469528}]F_{10}(T,\ldots,T)\bmod p&68&256&476348.
 \end{array}
\]
Any one residue proves the matching lower bound, so this degree is exact.
This is leading-coefficient propagation, not a random point evaluation.

The independent checker imports no builder. It independently parses all
22 source residuals; verifies the inverse row operations, 89 records,
and all 22,544 isolation interactions; checks 59 local circuit identities,
the additional Lucas factorization, and the eight-factor norm identity;
audits every gate and fixed binding; and certifies the exact degree.
Finite implementation tests cover 476 borrow transitions, 11,049 binary
block pairs, and 60 large-integer source assignments, including nonbinary
square bases 25 and 49. These finite checks supplement, rather than replace,
the unbounded mathematical arguments. Both ordinary Python and
\texttt{python3 -O independent\_check.py} pass.

Against the prior calibrated-mask ten-integer construction of degree
$21\,111\,400$ and 659 gates, the new point has $59.88\%$ lower degree and
126 fewer gates ($19.12\%$). Against the incoming 582-gate homogeneous
source point of degree $50\,822\,936$, it has 49 fewer gates and more than
sixfold lower degree. These are explicit pinned comparisons, not claims
of global optimality or superiority over every variable/domain tradeoff.

The inherited exact-source pin is
\texttt{124c07c0b77a};
\texttt{SOURCE\_PINS.json} preserves all 12 source identities.
Primary arithmetic dependencies are Z.-W. Sun,
\emph{Further results on Hilbert's Tenth Problem}, Sci. China Math. 64
(2021), 281--306, \url{https://arxiv.org/html/1704.03504v7}; and
\emph{Ten unknowns for Hilbert's tenth problem over the integers},
2026 preprint, \url{https://arxiv.org/html/2609.01594v2}.
The complete inherited source, coefficient, and finalizer proof records
are preserved under \texttt{upstream/}; their historical relative paths
do not supersede the active manifest.
